% Plan détaillé — version révisée après relecture.
% Compilation : pdflatex plan-rapport.tex
\documentclass[11pt,a4paper]{article}

\usepackage[T1]{fontenc}
\usepackage[utf8]{inputenc}
\usepackage{lmodern}
\usepackage[french]{babel}
\usepackage[a4paper,left=2.25cm,right=2.25cm,top=1.9cm,bottom=1.9cm]{geometry}
\usepackage{xcolor}
\usepackage{titlesec}
\usepackage{fancyhdr}
\usepackage{enumitem}
\usepackage{microtype}
\usepackage{textcomp}
\providecommand{\euro}{\texteuro}
\usepackage[hidelinks]{hyperref}

\definecolor{encre}{HTML}{1A1A19}
\definecolor{ocre}{HTML}{C79100}

\titleformat{\section}{\normalfont\large\bfseries\color{encre}}{}{0em}{}
\titlespacing*{\section}{0pt}{1.5ex plus 1ex minus .2ex}{0.6ex plus .2ex}

\pagestyle{fancy}
\fancyhf{}
\renewcommand{\headrulewidth}{0.4pt}
\fancyhead[L]{\small\itshape Un système fiscal pour le XXI\textsuperscript{e} siècle --- plan détaillé}
\fancyhead[R]{\small\thepage}

\setlength{\parindent}{0pt}
\setlength{\parskip}{0.20em}
\linespread{1.0}

\newenvironment{chapitres}[1]
{\begin{enumerate}[leftmargin=1.6em,itemsep=0.15em,parsep=0pt,topsep=0.20em,start=#1]}
{\end{enumerate}}

\begin{document}

\thispagestyle{empty}
{\footnotesize\color{ocre}\bfseries PLAN DÉTAILLÉ --- TROIS PARTIES\par}
\vspace{5pt}
{\raggedright\fontsize{18}{21}\selectfont\bfseries Un système fiscal\\ pour le XXI\textsuperscript{e} siècle\par}
\vspace{6pt}
{\color{encre}\hrule height 1.2pt}
\vspace{9pt}

\textbf{La thèse.} À niveau de recettes et de redistribution donnés, toutes les architectures
fiscales ne se valent pas. Il est possible de concevoir un système plus simple, plus neutre et moins
destructeur d'activité sans renoncer à ses objectifs de justice ou de financement. Cette thèse n'est
pas posée d'entrée : elle est la conclusion de la première partie, et c'est seulement une fois établie
que la deuxième peut mesurer l'écart français.

\section{Première partie --- Qu'est-ce qu'un bon impôt ?}

\begin{chapitres}{1}
\item \textbf{Le pacte fiscal : les principes qui fondent l'impôt.} Partir de la Déclaration de 1789
comme socle commun : la contribution est nécessaire, elle doit être répartie selon les facultés de
chacun, consentie et contrôlée. De là quelques exigences simples : l'impôt doit être légitime,
proportionné, compréhensible, et ne pas imposer de prélèvement ou de coût qui ne soit justifié par un
objectif collectif --- ce qui inclut les frais de gestion et le coût de conformité, qui font partie
du prix de l'impôt. C'est ici qu'on introduit l'idée qu'un impôt juste ne se réduit pas à son
barème : il doit aussi respecter les libertés, la capacité contributive et le consentement.
\item \textbf{Financer, orienter, corriger : le triangle de l'impôt.} Un système fiscal poursuit
plusieurs objectifs à la fois : financer l'action publique, modifier certains comportements lorsque
c'est justifié, corriger certaines inégalités ou défaillances. Ces objectifs entrent en tension.
L'intuition de Tinbergen la règle : lorsqu'on poursuit plusieurs objectifs, il est illusoire de
demander à un seul instrument de tous les atteindre correctement. D'où une idée qui commande tout le
rapport : un impôt isolé n'a pas à être parfait sur tous les critères, c'est le système fiscal et
social dans son ensemble qui doit être cohérent. Atkinson-Stiglitz le dit dans l'autre sens :
lorsqu'un impôt sur le revenu non linéaire existe, différencier les taux sur la consommation
n'ajoute rien.
\item \textbf{Prélever mieux : ce que nous apprend l'économie de l'impôt.} Les grandes intuitions de
la taxation optimale, expliquées sans technicité excessive, et chacune portée par un cas concret.
Pourquoi toute taxation peut créer une perte sèche --- un vendeur accepterait 200\,000~\euro{}, un
acheteur en offrirait 210\,000, et avec 8~\% de droits de mutation entre les deux aucun prix ne
convient : personne n'a payé l'impôt et les 10\,000~\euro{} de gain sont perdus pour tout le monde.
Pourquoi les taux marginaux commandent les comportements, puisque personne ne décide en fonction de
son taux moyen. Pourquoi une assiette réagit quand on la taxe, et jusqu'où va l'intuition de Laffer :
un sommet existe, mais sa position dépend de chaque impôt et de la facilité d'y échapper. Pourquoi
l'incidence réelle ne correspond pas toujours au payeur légal. Pourquoi, enfin, des assiettes larges
et des taux bas permettent de lever une recette donnée avec moins de distorsions et moins
d'arbitrages. C'est aussi ici qu'on explique pourquoi la redistribution doit être appréciée sur
l'ensemble impôts-transferts plutôt que taxe par taxe, et qu'on introduit le classement des impôts
par effet sur la croissance établi à l'OCDE, avec sa réserve : ses auteurs jugent eux-mêmes excessive
l'ampleur des effets estimés.
\end{chapitres}

La partie se referme sur la thèse, qui en découle : à niveau de recettes et de redistribution donnés,
toutes les architectures fiscales ne se valent pas. Si ces principes sont relativement simples,
comment la France s'en est-elle à ce point éloignée ?

\section{Deuxième partie --- Comment la France a perdu le fil}

Partie narrative, lisible d'une traite. La comparaison internationale n'a pas de chapitre autonome :
elle sert de test à chaque anomalie --- est-elle réellement française, et d'autres pays
atteignent-ils le même degré de redistribution avec une architecture différente ?

\begin{chapitres}{4}
\item \textbf{Un système que personne n'a vraiment dessiné.} L'économie politique de l'impôt, et
quelques exemples historiques frappants : la contribution des portes et fenêtres, créée en 1798 et
supprimée en 1926, choisie parce qu'elle était facile à constater, a muré les façades pendant un
siècle. Puis le récit français par strates. 1945 : asseoir l'assurance sociale sur le salaire
convenait au plein emploi et à une population active en croissance, et lie pour un demi-siècle le
financement de la protection sociale au coût du travail. 1991 : la CSG est la tentative d'en sortir,
mais l'élargissement s'arrête à mi-chemin ; la CRDS créée en 1996 pour treize ans est toujours là.
Depuis 1993 : des allègements empilés, le CICE, sa bascule en allègements, chacun avec son seuil et
son plafond. Chaque époque répond à son problème, personne ne refond l'ensemble : à l'impôt s'ajoute
bientôt la niche, puis la compensation, puis l'aide.
\item \textbf{Beaucoup prélever, mais surtout mal prélever.} Photographie actuelle et comparaison
internationale. Le niveau global d'abord : la France dépense 57,0~\% du PIB contre 49,4~\% en moyenne
dans la zone euro et prélève 45,3~\% contre 40,8~\%, le solde étant emprunté (Eurostat, 2024) --- le
niveau de prélèvement est la contrepartie du niveau de dépense. Puis surtout la structure, qui est un
problème indépendant : on peut prélever 40~\% du PIB aussi mal que 45. Reprendre alors la grille de
la première partie et passer le système au crible --- assiettes étroites et taux élevés, taxation du
travail et de la production, niches et différences de traitement, taux marginaux, fiscalité du
capital et de la consommation, bases immobiles mal exploitées. C'est le chapitre du « où la France
viole-t-elle le plus clairement nos principes ? », et il est chiffré : sur l'euro suivant, un
célibataire sans enfant est prélevé à 64,6~\% autour de 67~\% du salaire moyen contre 58,2~\% au
salaire moyen ; 77~milliards sont dus même quand l'entreprise perd de l'argent ; le même rendement
d'épargne est imposé de 0 à 55~\% selon l'enveloppe ; et la France prélève 25,9 points de PIB dans
l'avant-dernière catégorie du classement de l'OCDE contre 17,5 en moyenne.
\item \textbf{Quand la complexité finit par miner l'efficacité et la justice.} On taxe puis on
subventionne : on prélève lourdement le travail puis on crée des allègements, on taxe la production
puis on verse des aides, on institue un impôt général puis un taux réduit, on crée une niche puis un
plafonnement de niches. Le système produit ses propres antidotes. Il crée des arbitrages et de
l'optimisation ; les taux affichés ne correspondent plus aux taux effectivement payés --- 474
dépenses fiscales pour 91,8~milliards, dont la Cour des comptes écrit que le coût n'est pas connu. La
redistribution devient illisible : impôts progressifs, taxes proportionnelles, cotisations,
exonérations, prestations, avantages liés à l'âge ou au statut, et plus personne ne sait dire qui
paie quoi, d'où des controverses qu'aucun chiffre ne tranche. Le consentement de l'article 14 s'use,
faute de pouvoir constater quoi que ce soit. Et surtout : pourquoi les mécanismes politiques
favorisent la conservation des niches. Les pertes d'une réforme sont concentrées et ses gains
diffus ; une niche a des bénéficiaires qui savent exactement ce qu'ils perdent, quand celui qui gagne
cent euros à une baisse générale de taux ignore d'où ils viennent ; et certains impôts survivent
parce qu'ils sont peu visibles ou faciles à collecter.
\end{chapitres}

Le problème français n'est donc pas seulement que nous prélevons beaucoup, ni que notre système est
compliqué : nous avons progressivement substitué l'accumulation de prélèvements, d'exceptions et de
compensations à une architecture fiscale cohérente. Si l'on cessait de corriger à la marge et qu'on
essayait enfin de redessiner, que ferait-on ?

\section{Troisième partie --- Quel système fiscal pour la France ?}

Le mouvement est le suivant : voilà le système qu'on peut construire, voilà les choix qui restent
ouverts à l'intérieur de ce système, voilà pourquoi il est difficile d'y arriver.

\begin{chapitres}{7}
\item \textbf{À quoi ressemblerait concrètement un système fiscal plus simple et plus efficace ?}
Passer des principes aux instruments : un impôt sur le revenu à assiette plus large et à taux plus
faibles ; beaucoup moins de niches et d'abattements, y compris ceux qui tiennent au statut, comme
l'abattement de 10~\% sur les pensions ; une TVA avec beaucoup moins de taux réduits ; une fiscalité
de l'épargne plus homogène ; une fiscalité pigouvienne assumée lorsqu'il s'agit de corriger une
externalité, à commencer par la taxe carbone ; moins de taxation des intrants et de la production,
selon le principe « zéro aide contre zéro impôt de production ». L'idée est de montrer que le système
cible n'est pas abstrait : on peut réellement lever les recettes autrement. Le chapitre le montre par
un tableau en deux colonnes --- ce que rapporte chaque assiette aujourd'hui, ce qu'elle rapporterait
après réforme --- et par l'effet sur quelques cas types et sur les déciles publiés.
\item \textbf{Les grands arbitrages : où déplacer la charge fiscale ?} Ici, on traite les vrais
débats. Faut-il déplacer une partie de la taxation du travail vers la consommation, donc le débat sur
la TVA sociale, avec les précédents allemand de 2007 et danois et l'échec français de 2012 ? Quelle
place donner à la taxation des transmissions et de l'héritage --- présentée sous forme d'encadré
contradictoire : pour, c'est le revenu le moins lié à l'effort et le plus concentré ; contre, c'est
l'impôt le plus rejeté, qui frappe une épargne déjà imposée et dont le rendement dépend de la
capacité à retenir des assiettes mobiles ? Comment simplifier la fiscalité de l'épargne et arbitrer
entre rendement normal, rente et prise de risque ? Le chapitre assume qu'une fois les grands
principes posés, certains choix restent discutables et dépendent aussi de préférences collectives.
\item \textbf{Pourquoi passer au système cible est si difficile.} La suppression des niches fait des
perdants très visibles, quand les baisses générales de taux font des gagnants diffus. Il faut
préserver le rendement et la redistribution, et compenser les ménages réellement perdants. Certaines
baisses d'impôt, notamment pour les entreprises, sont politiquement difficiles. La transition doit
être progressive, et le risque principal est que les hausses soient votées avant les compensations
--- d'où la nécessité de réformer par paquets cohérents, chaque suppression votée dans le même texte
que sa compensation.
\end{chapitres}

\section{Ce qui reste à trancher}

\begin{enumerate}[leftmargin=1.6em,itemsep=0.15em,parsep=0pt,topsep=0.20em]
\item \textbf{La cible de dépense.} La moyenne de la zone euro est à 49,4~\% du PIB, 7,6 points sous
la France ; 52~\% est à mi-chemin, environ 147~milliards. Il faut choisir : 52~\% n'est pas « la
moyenne de la zone euro ».
\item \textbf{Le tableau du chapitre~7.} C'est lui qui montre que les recettes peuvent être levées
autrement : si les comptes ne tombent pas, certaines propositions changeront. Reste à décider qui le
construit et sur quelles sources publiques.
\item \textbf{L'assiette des transmissions.} Le niveau reste un arbitrage ouvert, mais les deux
positions s'accordent sur le caractère indéfendable des exceptions actuelles. Le rapport le
dit-il, ou l'encadré se clôt-il sans recommandation ?
\end{enumerate}

\end{document}
